\documentclass[12pt]{article}
\usepackage{amsmath}
\usepackage{amssymb}
\usepackage{amsthm}
\usepackage{parskip}
\begin{document}
\noindent \textbf{Problem 2.}
\bigbreak
\noindent \textbf{(a)}
\smallbreak
the truth table shows that the $(P \land \neg Q) \lor (P \land Q)$ is equivalent to $P$.
$$
\begin{array}{cc|cccc}
\hline
P & Q &  \neg Q  &  P \land \neg Q  &  P \land Q  &  (P \land \neg Q) \lor (P \land Q) \\
\hline\hline
T & T &   F   &    F      &    T    &            T         \\
T & F &   T   &    T      &    F    &            T         \\
F & T &   F   &    F      &    F    &            F         \\
F & F &   T   &    F      &    F    &            F         \\
\hline
\end{array}
$$
\smallbreak
\bigbreak
\noindent \textbf{(b)}
\smallbreak
$(A - B)$ is equivalent to $A \land \neg B$
\smallbreak
the truth table about $(A - B) \lor (A \land B)$ is
$$
\begin{array}{cc|cccc}
\hline
A & B & \neg B &  A \land \neg B  &  A \land B  & (A - B) \lor (A \land B)\\
\hline\hline
T & T &  F  &    F      &    T    &         T\\
T & F &  T  &    T      &    F    &         T\\
F & T &  F  &    F      &    F    &         F\\
F & F &  T  &    F      &    F    &         F\\
\hline
\end{array}
$$
\smallbreak
we can see the $A$ is equivalent to $(A - B) \lor (A \land B)$
\smallbreak
so $x \in A \iff x \in ((A - B) \lor (A \land B))$
\bigbreak
\bigbreak

\noindent \textbf{Problem 3.}
\bigbreak
\noindent \textbf{(a)}
sixish($n$) := $\exists m \exists k.((6 \cdot 10^k + m = n) \land m < 10^k )$
\smallbreak
\noindent \textbf{(b)}
$\exists m \exists n \forall k.Q(m) \land Q(n) \land k \neq m \land k \neq n \land \neg Q(k)$
\smallbreak
\noindent \textbf{(c)}
$\forall k \exists n \exists m.(isEven(k) \land k > 2 ) \implies (m + n = k) \land isPrime(n) \land isPrime(m)$
\smallbreak
\noindent \textbf{(d)}
$\forall n.n > 100 \implies $sixish(n)
\smallbreak
\noindent \textbf{(e)}
it is not equivalent because when n = 1, R(n) should be True but we can't find any b $\in \mathbb{N}$ to make 1 = 2b + 5
\bigbreak
\bigbreak

\noindent \textbf{Problem 4.}
\bigbreak
\noindent \textbf{(a)}

We use induction.

\begin{proof}
    Let $P(n)$ be the statement that if $Q(x_1 \cdot x_2 \cdots x_n)$ is true, then at least one of $Q(x_1)$, $Q(x_2)$, $\dots$, $Q(x_n)$ must be true.

    \textbf{Base case:} For $n = 1$, the statement holds trivially.For $n = 2$, if $Q(x_1 \cdot X_2)$ is true, then at least one of $Q(x_1)$ or $Q(x_2)$ must be true which holds by the given premise.

    \textbf{Inductive hypothesis:} Assume that $P(n)$ is true for some $n \ge 1$. That is, if $Q(x_1 \cdot x_2 \cdots x_n)$ is true, then at least one of $Q(x_1), \dots, Q(x_n)$ must be true.

    \textbf{Inductive step}: We assume that for $n \ge 1$, P($n$) is true. We have to prove $P(n + 1)$ is true, which means we have to prove if $Q(x_1 \cdot x_2 \cdots x_{n + 1})$ is true,then at least one of $Q(x_1)$ $\ldots$ $Q(x_{n + 1})$ must be true. Since the $x_1 \cdot x_{n + 1}$ can be written as $(x_1 \cdot x_n) \cdot x_{n + 1}$, so let $a$ be $x_1 \cdot x_2 \cdots x_n$ and b be $x_{n + 1}$, if $Q(a \cdot b)$ is true then one of $Q(a)$ or $Q(b)$ must be true.

    \begin{itemize}
        \item Case 1: if $Q(a)$ is true, then $Q(x_1 \cdots x_n)$ is true by the inductive hypothesis
        \item Cast 2: if $Q(b)$ is true, then $Q(x_{n + 1}$) is true which proved the P($n + 1$) is true directly.
    \end{itemize}

    So we proved $P(n)$ is true for all $n \ge 1$.

\end{proof}

\end{document}
